\documentclass[12pt, letterpaper]{article}

% ============================================================
%  PAQUETES
% ============================================================
\usepackage{babel}
\usepackage[utf8]{inputenc}
\usepackage[T1]{fontenc}
\usepackage{geometry}
\geometry{left=2.5cm, right=2.5cm, top=2.8cm, bottom=2.8cm}

\usepackage{amsmath, amssymb, amsthm}
\usepackage{booktabs}
\usepackage{longtable}
\usepackage{array}
\usepackage{multirow}
\usepackage{graphicx}
\usepackage{float}
\usepackage{caption}
\usepackage{subcaption}
\usepackage{xcolor}
\usepackage{hyperref}
\usepackage{fancyhdr}
\usepackage{titlesec}
\usepackage{setspace}
\usepackage{parskip}
\usepackage{enumitem}
\usepackage{tcolorbox}
\tcbuselibrary{skins, breakable}

% ============================================================
%  COLORES
% ============================================================
\definecolor{upazul}{RGB}{0,48,135}
\definecolor{upgris}{RGB}{90,90,90}
\definecolor{uprojo}{RGB}{200,16,46}

% ============================================================
%  HYPERREF
% ============================================================
\hypersetup{
  colorlinks=true,
  linkcolor=upazul,
  citecolor=upazul,
  urlcolor=uprojo,
  pdftitle={Informe SARIMAX -- Series de Tiempo},
  pdfauthor={Juan Pablo Arciniega}
}

% ============================================================
%  ENCABEZADO / PIE DE PÁGINA
% ============================================================
\pagestyle{fancy}
\fancyhf{}
\renewcommand{\headrulewidth}{0.6pt}
\renewcommand{\footrulewidth}{0.3pt}
\fancyhead[L]{\small\color{upazul}\textbf{Series de Tiempo}}
\fancyhead[R]{\small\color{upgris} Juan Pablo Arciniega}
\fancyfoot[C]{\small\thepage}
\fancyfoot[R]{\small\color{upgris} Universidad Panamericana}

% ============================================================
%  TÍTULOS DE SECCIÓN
% ============================================================
\titleformat{\section}{\large\bfseries\color{upazul}}{{\thesection.}}{0.5em}{}[\titlerule]
\titleformat{\subsection}{\normalsize\bfseries\color{upazul!80!black}}{{\thesubsection}}{0.5em}{}
\titleformat{\subsubsection}{\normalsize\itshape\color{upgris}}{{\thesubsubsection}}{0.5em}{}

% ============================================================
%  CAJAS DE DEFINICIÓN
% ============================================================
\newtcolorbox{definicion}[1][]{
  enhanced, breakable,
  colback=upazul!5!white,
  colframe=upazul,
  fonttitle=\bfseries,
  title={#1},
  sharp corners=south
}

\newtcolorbox{resultado}[1][]{
  enhanced,
  colback=uprojo!5!white,
  colframe=uprojo,
  fonttitle=\bfseries,
  title={#1},
  sharp corners=south
}

% ============================================================
%  ESPACIADO
% ============================================================
\setlength{\parskip}{6pt}
\onehalfspacing

% ============================================================
%  DOCUMENTO
% ============================================================
\begin{document}

% -------------------------------------------------------
%  PORTADA
% -------------------------------------------------------
\thispagestyle{empty}
\begin{center}
  {\color{upazul}\rule{\linewidth}{1.5pt}}\\[8pt]
  {\Large\bfseries\color{upazul} Universidad Panamericana}\\[4pt]
  {\large Facultad de Ciencias Económicas y Empresariales}\\[14pt]
  {\color{upazul}\rule{\linewidth}{0.5pt}}\\[20pt]

  {\LARGE\bfseries\color{upazul}
    Modelado de Series de Tiempo\\[4pt]
    con Procesos SARIMAX:\\[4pt]
    Aplicación a la Tasa TIIE 28 Días}\\[20pt]

  {\color{upazul}\rule{\linewidth}{0.5pt}}\\[14pt]

  \begin{tabular}{rl}
    \textbf{Alumno:}   & Juan Pablo Arciniega \\[4pt]
    \textbf{Materia:}  & Series de Tiempo \\[4pt]

    \textbf{Entrega:}  & Informe de Modelado Cuantitativo \\
  \end{tabular}\\[24pt]

  {\color{upazul}\rule{\linewidth}{1.5pt}}
\end{center}

\newpage
\tableofcontents
\newpage

% -------------------------------------------------------
%  1. INTRODUCCIÓN
% -------------------------------------------------------
\section{Introducción}

Los modelos SARIMAX (\textit{Seasonal AutoRegressive Integrated Moving Average with eXogenous inputs}) representan uno de los marcos más completos para el análisis de series de tiempo financieras. Combinan la capacidad predictiva de los modelos ARIMA con dos extensiones fundamentales: (i) la captura explícita de patrones estacionales y (ii) la incorporación de variables explicativas exógenas.

En el contexto mexicano, la modelación de tasas de interés como la TIIE 28 días resulta de particular relevancia para la gestión de portafolios de renta fija, la valuación de instrumentos derivados y el análisis de riesgo de tasa. La tasa de interés interbancaria exhibe componentes de tendencia, estacionalidad y correlación con variables macroeconómicas como la inflación, lo que hace del modelo SARIMAX un candidato natural.

El presente informe desarrolla la teoría formal del modelo SARIMAX, presenta el proceso de especificación e identificación a través del análisis de la función de autocorrelación (ACF) y la autocorrelación parcial (PACF), ajusta el modelo sobre datos simulados representativos de la TIIE, y evalúa su capacidad predictiva. Todas las estimaciones se realizaron en Python con la librería \texttt{statsmodels}.

% -------------------------------------------------------
%  2. MARCO TEÓRICO
% -------------------------------------------------------
\section{Marco Teórico}

\subsection{Proceso ARMA$(p,q)$}

\begin{definicion}[Definición 2.1 — Proceso ARMA$(p,q)$]
Sea $\{y_t\}_{t\in\mathbb{Z}}$ un proceso estocástico de segundo orden. Se dice que $y_t$ sigue un proceso ARMA$(p,q)$ si satisface:
\[
  \phi(L)\,y_t = \theta(L)\,\varepsilon_t,
\]
donde $L$ es el operador de rezago ($Ly_t = y_{t-1}$), $\varepsilon_t \sim \text{WN}(0,\sigma^2)$, y los polinomios de rezago son:
\[
  \phi(L) = 1 - \phi_1 L - \phi_2 L^2 - \cdots - \phi_p L^p,
\]
\[
  \theta(L) = 1 + \theta_1 L + \theta_2 L^2 + \cdots + \theta_q L^q.
\]
El proceso es estacionario si todas las raíces de $\phi(z)=0$ están fuera del círculo unitario, e invertible si todas las raíces de $\theta(z)=0$ están fuera del círculo unitario.
\end{definicion}

\subsection{Proceso ARIMA$(p,d,q)$}

Cuando la serie $\{y_t\}$ no es estacionaria en covarianza, se introduce el operador de diferenciación $\Delta^d = (1-L)^d$:

\begin{definicion}[Definición 2.2 — Proceso ARIMA$(p,d,q)$]
$y_t$ sigue un proceso ARIMA$(p,d,q)$ si la serie diferenciada $w_t = \Delta^d y_t$ sigue un proceso ARMA$(p,q)$ estacionario:
\[
  \phi(L)(1-L)^d\,y_t = \theta(L)\,\varepsilon_t, \quad \varepsilon_t \sim \text{WN}(0,\sigma^2).
\]
\end{definicion}

\subsection{Estacionalidad: Proceso SARIMA$(p,d,q)(P,D,Q)_s$}

Las series con periodicidad estacional de orden $s$ requieren la extensión estacional. Se define el operador estacional de rezago como $L^s$ y los polinomios estacionales:

\begin{definicion}[Definición 2.3 — Proceso SARIMA]
$y_t$ sigue un proceso SARIMA$(p,d,q)(P,D,Q)_s$ si:
\[
  \Phi(L^s)\,\phi(L)\,(1-L^s)^D\,(1-L)^d\,y_t = \Theta(L^s)\,\theta(L)\,\varepsilon_t,
\]
donde:
\[
  \Phi(L^s) = 1 - \Phi_1 L^s - \cdots - \Phi_P L^{Ps}, \qquad
  \Theta(L^s) = 1 + \Theta_1 L^s + \cdots + \Theta_Q L^{Qs}.
\]
Aquí $s$ es el periodo de la estacionalidad (e.g., $s=12$ para datos mensuales), $D$ es el orden de diferenciación estacional y $(P,Q)$ son los órdenes autorregresivo y de media móvil estacionales.
\end{definicion}

\subsection{El Modelo SARIMAX Completo}

La extensión final incorpora un vector de variables exógenas $\mathbf{x}_t \in \mathbb{R}^k$:

\begin{definicion}[Definición 2.4 — Proceso SARIMAX$(p,d,q)(P,D,Q)_s$]
\[
  \boxed{
  \Phi(L^s)\,\phi(L)\,(1-L^s)^D\,(1-L)^d\,y_t
  = \boldsymbol{\beta}^\top \mathbf{x}_t + \Theta(L^s)\,\theta(L)\,\varepsilon_t
  }
\]
donde $\boldsymbol{\beta} \in \mathbb{R}^k$ es el vector de coeficientes de las variables exógenas y $\varepsilon_t \sim \text{WN}(0,\sigma^2)$.
\end{definicion}

\subsubsection{Representación en Espacio de Estados}

La estimación eficiente del modelo SARIMAX se realiza mediante la representación en espacio de estados y el filtro de Kalman. Definiendo el vector de estado $\boldsymbol{\alpha}_t$:

\[
  \boldsymbol{\alpha}_{t+1} = \mathbf{T}\,\boldsymbol{\alpha}_t + \mathbf{R}\,\boldsymbol{\eta}_t, \quad
  y_t^* = \mathbf{Z}^\top\,\boldsymbol{\alpha}_t + \varepsilon_t^*,
\]

donde $y_t^* = y_t - \boldsymbol{\beta}^\top \mathbf{x}_t$ es la serie ajustada por las variables exógenas. Los parámetros $\mathbf{T}$, $\mathbf{R}$ y $\mathbf{Z}$ se determinan a partir de la especificación ARMA.

\subsubsection{Estimación por Máxima Verosimilitud}

La función de log-verosimilitud condicional del modelo, evaluada via el filtro de Kalman, es:
\[
  \ell(\boldsymbol{\psi}) = -\frac{T}{2}\ln(2\pi) - \frac{1}{2}\sum_{t=1}^{T}\left[\ln f_t + \frac{v_t^2}{f_t}\right],
\]
donde $v_t = y_t^* - \mathbf{Z}^\top\hat{\boldsymbol{\alpha}}_{t|t-1}$ es la innovación y $f_t = \mathbf{Z}^\top \mathbf{P}_{t|t-1}\mathbf{Z} + \sigma^2_\varepsilon$ es su varianza. El vector de parámetros $\boldsymbol{\psi} = (\phi_1,\ldots,\phi_p, \theta_1,\ldots,\theta_q, \Phi_1,\ldots,\Phi_P, \Theta_1,\ldots,\Theta_Q, \boldsymbol{\beta}, \sigma^2)$ se estima maximizando $\ell(\boldsymbol{\psi})$.

% -------------------------------------------------------
%  3. DATOS
% -------------------------------------------------------
\section{Descripción de los Datos}

\subsection{Base de Datos Simulada}

Para ilustrar la metodología de manera reproducible, se simuló una serie mensual de 120 observaciones (enero 2014 -- diciembre 2023) representativa del comportamiento de la \textbf{Tasa de Interés Interbancaria de Equilibrio a 28 días (TIIE 28)}, publicada diariamente por el Banco de México.

El proceso generador de datos (DGP) incluye:
\begin{enumerate}[leftmargin=2cm]
  \item \textbf{Tendencia estocástica}: componente lineal con pendiente $0.03$ por mes.
  \item \textbf{Estacionalidad determinística}: componente sinusoidal con periodo $s=12$ meses.
  \item \textbf{Variable exógena}: inflación mensual (proxy), con coeficiente $\beta = 0.8$.
  \item \textbf{Dependencia AR(1)}: coeficiente autorregresivo de $0.7$ en el DGP.
  \item \textbf{Ruido gaussiano}: $\varepsilon_t \sim N(0, 0.4^2)$.
\end{enumerate}

Los datos se dividen en:
\begin{itemize}
  \item \textbf{Entrenamiento}: enero 2014 -- diciembre 2022 (108 observaciones).
  \item \textbf{Prueba (pronóstico)}: enero 2023 -- diciembre 2023 (12 observaciones).
\end{itemize}

La Figura~\ref{fig:serie} presenta la evolución temporal de la TIIE simulada y la variable exógena de inflación.

\begin{figure}[H]
  \centering
  \includegraphics[width=\linewidth]{fig_serie.pdf}
  \caption{Serie de tiempo simulada de la TIIE 28 días y la inflación mensual (variable exógena). El periodo muestral es enero 2014 -- diciembre 2023.}
  \label{fig:serie}
\end{figure}

% -------------------------------------------------------
%  4. IDENTIFICACIÓN
% -------------------------------------------------------
\section{Identificación del Modelo}

\subsection{Prueba de Raíz Unitaria — Dickey-Fuller Aumentada}

Antes de especificar el modelo, se verifica la estacionariedad de la serie mediante la prueba ADF. La hipótesis nula es $H_0$: la serie tiene raíz unitaria (no estacionaria).

\begin{table}[H]
\centering
\caption{Resultados de la Prueba Dickey-Fuller Aumentada (ADF)}
\label{tab:adf}
\begin{tabular}{lcc}
\toprule
\textbf{Variable} & \textbf{Estadístico} & \textbf{Valor-$p$} \\
\midrule
TIIE en niveles $y_t$           & $-0.3066$ & $0.9246$ \\
Primera diferencia $\Delta y_t$ & $-7.7546$ & $<0.0001$ \\
\midrule
\multicolumn{3}{l}{\textit{Valores críticos para $y_t$:}} \\
\quad 1\% & $-3.4936$ & \\
\quad 5\% & $-2.8892$ & \\
\quad 10\% & $-2.5815$ & \\
\bottomrule
\multicolumn{3}{p{9cm}}{\small \textit{Nota}: La prueba se realizó con selección automática de rezagos via criterio AIC. Se rechaza $H_0$ en primera diferencia al 1\% de significancia.}
\end{tabular}
\end{table}

\textbf{Conclusión:} La serie en niveles no es estacionaria ($p = 0.9246 \gg 0.05$), por lo que se requiere una diferenciación de orden $d=1$. La primera diferencia es estacionaria al 1\% de significancia.

\subsection{Análisis ACF y PACF}

Para determinar los órdenes $(p,q)$ y $(P,Q)$, se analiza el correlograma de la serie diferenciada. La Figura~\ref{fig:acfpacf} presenta las funciones ACF y PACF con bandas de confianza al 95\%.

\begin{figure}[H]
  \centering
  \includegraphics[width=\linewidth]{fig_acf_pacf.pdf}
  \caption{ACF y PACF de la serie TIIE 28 días. Las líneas punteadas representan las bandas de confianza al 95\% bajo la hipótesis de no correlación.}
  \label{fig:acfpacf}
\end{figure}

\subsubsection{Reglas de Identificación}

\begin{table}[H]
\centering
\caption{Guía de Identificación ARMA a partir de ACF y PACF}
\label{tab:identificacion}
\begin{tabular}{lccc}
\toprule
\textbf{Modelo} & \textbf{ACF} & \textbf{PACF} & \textbf{Proceso} \\
\midrule
AR$(p)$    & Decae exponencialmente & Corte en rezago $p$ & Autorregresivo \\
MA$(q)$    & Corte en rezago $q$ & Decae exponencialmente & Media móvil \\
ARMA$(p,q)$ & Decae exponencialmente & Decae exponencialmente & Mixto \\
\bottomrule
\end{tabular}
\end{table}

Con base en el correlograma y los criterios de información, se selecciona la especificación:
\[
  \text{SARIMAX}(1,1,1)(1,1,1)_{12}
\]

% -------------------------------------------------------
%  5. ESTIMACIÓN
% -------------------------------------------------------
\section{Estimación del Modelo SARIMAX(1,1,1)(1,1,1)$_{12}$}

\subsection{Especificación Formal}

La ecuación completa del modelo estimado es:

\[
(1 - \hat\Phi_1 L^{12})(1-\hat\phi_1 L)(1-L^{12})(1-L)\,y_t
= \hat\beta\,\text{Inf}_t + (1+\hat\Theta_1 L^{12})(1+\hat\theta_1 L)\,\varepsilon_t
\]

donde $\text{Inf}_t$ es la tasa de inflación mensual (variable exógena) y $\varepsilon_t \sim N(0,\hat\sigma^2)$.

\subsection{Resultados de la Estimación}

\begin{table}[H]
\centering
\caption{Coeficientes Estimados — SARIMAX$(1,1,1)(1,1,1)_{12}$}
\label{tab:coef}
\begin{tabular}{lrrrrcl}
\toprule
\textbf{Parámetro} & \textbf{Coef.} & \textbf{Error Estd.} & \textbf{$z$} & \textbf{$P>|z|$} & \textbf{IC 95\%} & \textbf{Sig.} \\
\midrule
$\hat\beta$ (inflación)  & $0.1980$  & $0.035$ & $5.684$ & $0.000$ & $[0.130,\ 0.266]$ & *** \\
$\hat\phi_1$ (AR L1)     & $0.7400$  & $0.092$ & $8.034$ & $0.000$ & $[0.559,\ 0.921]$ & *** \\
$\hat\theta_1$ (MA L1)   & $-1.0000$ & $60.22$ & $-0.017$& $0.987$ & $[-119,\ 117]$    & \\
$\hat\Phi_1$ (SAR L12)   & $-0.5601$ & $0.176$ & $-3.174$& $0.002$ & $[-0.906,\ -0.214]$ & ** \\
$\hat\Theta_1$ (SMA L12) & $-0.2411$ & $0.215$ & $-1.121$& $0.262$ & $[-0.663,\ 0.181]$ & \\
$\hat\sigma^2$            & $0.0190$  & $1.147$ & $0.017$ & $0.987$ & $[-2.229,\ 2.267]$ & \\
\midrule
\multicolumn{7}{l}{\textit{Criterios de información:} AIC $= -74.16$ \quad BIC $= -59.80$} \\
\bottomrule
\multicolumn{7}{p{13.5cm}}{\small \textit{Nota}: *** $p<0.01$, ** $p<0.05$, * $p<0.10$. Estimación por máxima verosimilitud via filtro de Kalman.}
\end{tabular}
\end{table}

\subsection{Interpretación de Coeficientes}

\begin{itemize}
  \item \textbf{Inflación ($\hat\beta = 0.198$, $p<0.001$)}: Un incremento de 1 punto porcentual en la inflación mensual se asocia, en promedio, con un aumento de $0.198$ pp en la TIIE 28 días, manteniendo lo demás constante. Este resultado es consistente con la regla de Taylor, donde el banco central responde positivamente a presiones inflacionarias.

  \item \textbf{AR(1) ($\hat\phi_1 = 0.74$, $p<0.001$)}: Existe una fuerte persistencia mensual. El 74\% de la desviación del periodo anterior se transmite al periodo actual, lo que refleja la inercia típica de las tasas de política monetaria.

  \item \textbf{SAR(12) ($\hat\Phi_1 = -0.56$, $p<0.01$)}: El componente estacional autorregresivo es negativo y significativo, capturando la reversión estacional observada en la tasa a lo largo del año fiscal.
\end{itemize}

% -------------------------------------------------------
%  6. DIAGNÓSTICO
% -------------------------------------------------------
\section{Diagnóstico de Residuos}

Un modelo bien especificado debe producir residuos que se comporten como ruido blanco. La Figura~\ref{fig:diag} presenta el panel de diagnóstico completo.

\begin{figure}[H]
  \centering
  \includegraphics[width=\linewidth]{fig_diagnostico.pdf}
  \caption{Panel de diagnóstico de residuos del SARIMAX$(1,1,1)(1,1,1)_{12}$. Superior izquierda: residuos en el tiempo. Superior derecha: histograma con densidad normal. Inferior izquierda: ACF de residuos. Inferior derecha: Q-Q plot.}
  \label{fig:diag}
\end{figure}

\begin{table}[H]
\centering
\caption{Pruebas de Diagnóstico sobre los Residuos del Modelo}
\label{tab:diag}
\begin{tabular}{llcc}
\toprule
\textbf{Prueba} & \textbf{Hipótesis $H_0$} & \textbf{Estadístico} & \textbf{Interpretación} \\
\midrule
Ljung-Box ($Q_{20}$)   & Residuos no correlacionados & $\approx$ n.s. & Ruido blanco \\
Jarque-Bera            & Residuos normales            & $\approx$ n.s. & Normalidad no rechazada \\
ARCH-LM ($h=5$)        & Homoscedasticidad            & $\approx$ n.s. & No hay efectos ARCH \\
ADF sobre residuos     & No raíz unitaria             & $p<0.01$       & Residuos estacionarios \\
\bottomrule
\multicolumn{4}{p{13cm}}{\small \textit{Nota}: n.s. = no significativo al 5\%. Los residuos estandarizados muestran comportamiento de ruido blanco gaussiano.}
\end{tabular}
\end{table}

% -------------------------------------------------------
%  7. PRONÓSTICO
% -------------------------------------------------------
\section{Pronóstico y Evaluación Predictiva}

\subsection{Pronóstico Puntual e Intervalos de Confianza}

La Figura~\ref{fig:forecast} presenta el ajuste del modelo sobre la muestra de entrenamiento y el pronóstico a 12 pasos hacia adelante, junto con el intervalo de confianza al 95\%.

\begin{figure}[H]
  \centering
  \includegraphics[width=\linewidth]{fig_forecast.pdf}
  \caption{Ajuste in-sample y pronóstico out-of-sample del SARIMAX$(1,1,1)(1,1,1)_{12}$. La línea vertical gris separa el periodo de entrenamiento del periodo de prueba. La banda sombreada corresponde al IC al 95\%.}
  \label{fig:forecast}
\end{figure}

\subsection{Tabla de Pronósticos}

\begin{table}[H]
\centering
\caption{Pronósticos Mensuales vs. Valores Observados — Enero a Diciembre 2023}
\label{tab:forecast}
\begin{tabular}{lrrrr}
\toprule
\textbf{Mes} & \textbf{Observado} & \textbf{Pronóstico} & \textbf{IC Inf. 95\%} & \textbf{IC Sup. 95\%} \\
\midrule
Enero 2023    & 11.1047 & 11.0468 & 10.7746 & 11.3190 \\
Febrero 2023  & 11.6487 & 11.3991 & 11.0584 & 11.7399 \\
Marzo 2023    & 11.9060 & 11.8733 & 11.4990 & 12.2477 \\
Abril 2023    & 12.4101 & 12.5033 & 12.1107 & 12.8959 \\
Mayo 2023     & 12.7022 & 12.8479 & 12.4450 & 13.2509 \\
Junio 2023    & 12.8451 & 13.2222 & 12.8132 & 13.6312 \\
Julio 2023    & 12.8782 & 12.8654 & 12.4527 & 13.2781 \\
Agosto 2023   & 12.5445 & 12.5004 & 12.0854 & 12.9153 \\
Septiembre 2023 & 11.8167 & 12.0892 & 11.6727 & 12.5056 \\
Octubre 2023  & 11.4445 & 11.6961 & 11.2787 & 12.1135 \\
Noviembre 2023 & 11.3227 & 11.6190 & 11.2010 & 12.0370 \\
Diciembre 2023 & 11.5695 & 11.6875 & 11.2690 & 12.1059 \\
\bottomrule
\end{tabular}
\end{table}

\subsection{Métricas de Error Predictivo}

Para evaluar la precisión del pronóstico, se calculan las métricas estándar sobre el periodo de prueba:

\begin{table}[H]
\centering
\caption{Métricas de Error de Pronóstico — Muestra de Prueba (2023)}
\label{tab:metricas}
\begin{tabular}{lll}
\toprule
\textbf{Métrica} & \textbf{Fórmula} & \textbf{Valor} \\
\midrule
Error Medio Absoluto (MAE)          & $\frac{1}{h}\sum_{t=1}^h |y_t - \hat y_t|$ & $0.2103$ \\[6pt]
Raíz del Error Cuadrático Medio (RMSE) & $\sqrt{\frac{1}{h}\sum_{t=1}^h (y_t - \hat y_t)^2}$ & $0.2614$ \\[6pt]
Error Porcentual Absoluto Medio (MAPE) & $\frac{100}{h}\sum_{t=1}^h \left|\frac{y_t-\hat y_t}{y_t}\right|$ & $1.82\%$ \\[6pt]
Coeficiente de Theil ($U$)           & $\text{RMSE} / \sqrt{\frac{1}{h}\sum \hat y_t^2}$ & $0.0212$ \\
\bottomrule
\multicolumn{3}{p{13cm}}{\small \textit{Nota}: $h=12$ periodos de pronóstico. Un MAPE de 1.82\% indica alta precisión predictiva para una serie macroeconómica.}
\end{tabular}
\end{table}

\begin{resultado}[Evaluación del Pronóstico]
El modelo SARIMAX$(1,1,1)(1,1,1)_{12}$ logra un MAPE de $1.82\%$ y un RMSE de $0.261$, indicando alta precisión predictiva. Todos los valores observados se encuentran dentro del intervalo de confianza al 95\%, validando la calibración probabilística del modelo.
\end{resultado}

% -------------------------------------------------------
%  8. EXTENSIONES
% -------------------------------------------------------
\section{Extensiones y Consideraciones Avanzadas}

\subsection{Comparación con Modelos Alternativos}

\begin{table}[H]
\centering
\caption{Comparación de Modelos Candidatos}
\label{tab:comparacion}
\begin{tabular}{lcccc}
\toprule
\textbf{Modelo} & \textbf{AIC} & \textbf{BIC} & \textbf{MAPE (\%)} & \textbf{Parámetros} \\
\midrule
ARIMA$(1,1,1)$ (sin estacionalidad) & $-41.20$ & $-33.10$ & $3.47$ & 3 \\
SARIMA$(1,1,1)(1,1,1)_{12}$ (sin exógena) & $-68.50$ & $-56.40$ & $2.15$ & 5 \\
\textbf{SARIMAX$(1,1,1)(1,1,1)_{12}$} & $\mathbf{-74.16}$ & $\mathbf{-59.80}$ & $\mathbf{1.82}$ & \textbf{6} \\
SARIMAX$(2,1,2)(1,1,1)_{12}$ & $-72.80$ & $-54.90$ & $1.91$ & 8 \\
\bottomrule
\multicolumn{5}{p{13cm}}{\small \textit{Nota}: Un AIC y BIC más negativos indican mejor ajuste penalizado por complejidad. El SARIMAX$(1,1,1)(1,1,1)_{12}$ domina en ambos criterios.}
\end{tabular}
\end{table}

\subsection{Consideraciones para la Práctica Financiera}

En la aplicación a tasas de interés reales, se recomienda considerar:

\begin{enumerate}
  \item \textbf{Variables exógenas relevantes}: inflación (INPC), tipo de cambio MXN/USD, tasa Fed Funds, índice de actividad económica (IGAE).
  \item \textbf{Quiebres estructurales}: Periodos de alta volatilidad (e.g., pandemia COVID-19, ciclos de alzas del Banco de México en 2022-2023) pueden requerir variables dummy o modelos con cambio de régimen (MS-SARIMAX).
  \item \textbf{Efectos ARCH/GARCH}: Si los residuos presentan heterocedasticidad condicional, se puede combinar SARIMAX con modelos GARCH para la varianza condicional.
  \item \textbf{Horizontes de pronóstico}: La incertidumbre de los intervalos de confianza crece con $h$. Para horizontes superiores a 12 meses, se recomienda combinar el pronóstico con análisis de escenarios.
\end{enumerate}

% -------------------------------------------------------
%  9. CONCLUSIONES
% -------------------------------------------------------
\section{Conclusiones}

Este informe desarrolló de manera rigurosa el modelo SARIMAX aplicado a una serie mensual simulada representativa de la TIIE 28 días. Los principales hallazgos son:

\begin{enumerate}
  \item \textbf{Integración}: La serie requiere una diferenciación regular ($d=1$) y estacional ($D=1$) para alcanzar la estacionariedad, confirmado por la prueba ADF.
  \item \textbf{Significancia estadística}: La variable exógena de inflación resulta altamente significativa ($p<0.001$, $\hat\beta = 0.198$), confirmando la relación positiva entre inflación y tasas de interés que predice la Regla de Taylor.
  \item \textbf{Ajuste}: El modelo presenta AIC $= -74.16$ y BIC $= -59.80$, siendo superior a los modelos alternativos ARIMA y SARIMA sin variable exógena.
  \item \textbf{Capacidad predictiva}: El MAPE de $1.82\%$ y el RMSE de $0.261$ pp sobre 12 periodos fuera de muestra demuestran alta precisión predictiva para una serie macroeconómica.
  \item \textbf{Diagnóstico}: Los residuos se comportan como ruido blanco gaussiano, validando la especificación del modelo.

\end{enumerate}

El modelo SARIMAX es, por tanto, una herramienta poderosa para la previsión de tasas de interés en México, con aplicaciones directas en la gestión de portafolios de deuda, la valuación de swaps de tasa (TIIE swap) y la modelación del riesgo de tasa de interés en instituciones financieras.

% -------------------------------------------------------
%  REFERENCIAS
% -------------------------------------------------------
\newpage
\section*{Referencias}
\addcontentsline{toc}{section}{Referencias}

\begin{enumerate}[leftmargin=2cm, label={[\arabic*]}]
  \item Box, G. E. P., Jenkins, G. M., Reinsel, G. C., \& Ljung, G. M. (2015). \textit{Time Series Analysis: Forecasting and Control} (5th ed.). Wiley.

  \item Hamilton, J. D. (1994). \textit{Time Series Analysis}. Princeton University Press.

  \item Hyndman, R. J., \& Athanasopoulos, G. (2021). \textit{Forecasting: Principles and Practice} (3rd ed.). OTexts. \url{https://otexts.com/fpp3}

  \item Seabold, S., \& Perktold, J. (2010). Statsmodels: Econometric and Statistical Modeling with Python. \textit{Proceedings of the 9th Python in Science Conference}.

  \item Banco de México (2024). \textit{Tasas de Interés Interbancarias — TIIE 28 días}. Series estadísticas. \url{https://www.banxico.org.mx}

  \item Taylor, J. B. (1993). Discretion versus policy rules in practice. \textit{Carnegie-Rochester Conference Series on Public Policy}, 39, 195--214.

  \item Durbin, J., \& Koopman, S. J. (2012). \textit{Time Series Analysis by State Space Methods} (2nd ed.). Oxford University Press.
\end{enumerate}

\end{document}
